\clearpage
\begin{resumen}
\addcontentsline{toc}{chapter}{Resumen}
Este trabajo compara representaciones textuales léxicas y basadas en LLM para predecir resultados de licitaciones públicas a partir de documentos de licitación en un subconjunto procesado de 1.162 contrataciones de construcción vial de Paraguay. Bajo validación cruzada estratificada de cinco pliegues, TF--IDF con regresión logística produce el mejor ordenamiento, con ROC AUC medio de 0,814 y PR AUC de 0,899, mientras que un transformer liviano entre fragmentos, aplicado sobre embeddings congelados de fragmentos de un LLM, logra la mejor calibración, con un puntaje de Brier de 0,162, y el comportamiento más estable ante el ajuste del umbral. El estudio destaca que los documentos de contratación contienen señal predictiva temprana y que el ordenamiento, la calibración y la sensibilidad al umbral deben evaluarse por separado para el monitoreo. Artefactos: doi:10.5281/zenodo.20128768.

\textbf{Palabras clave:} contratación pública, gobierno electrónico, democracia electrónica, documentos de licitación, pliego de bases y condiciones, modelos grandes de lenguaje, predicción de resultados.
\end{resumen}
